\begin{abstract}
持久化的 AI 伴侣（persistent AI companions）是许多用户接触长上下文 Agent 的入口：
有人用数百轮对话养一个数字角色，然后它忘了自己的名字或口头禅，关掉会话它就不在了。
这种失效较少来自能力，更多来自注入面（injection surface）：记忆被当作当前事实来读，
风格信号软化了拒绝，人格约束逐渐衰减，而截断可能把安全指令挤出去。
我们提出 LINGSHU——一个把 prompt 组装视为控制面的确定性、可审计控制面。
六个插件共享同一个只读的宪法边界；一个五域合成器在它们之间裁决；
控制面在模型调用之前加入带抗饱和（anti-windup）的积分反馈、带人格锚点的相变前馈监控、
必要多样性降维，以及一道源自 STPA、在组装边界上结构性约束「什么能到达模型」的闸门。
恰好一个节点是 fail-closed——STPA 审计——因此违规会阻断注入而不是降级注入；
其余组件以 fail-open 方式降级并记录原因。契约可机器校验：14 个文件中的 892 条 Python 断言
与 396 条插件侧断言，副本之间的参数一致性按键比对。
一个混沌套件在模拟宿主与真实宿主上逐一驱动每个机制；v5.4 的同回合投递路径让相位状态
落在携带临界比率的那次请求里，并由 24 条独立回归断言钉住。一次真实宿主 smoke 运行显示
两层已经解耦：闸门以 soft 重定向拦下一次角色扮演越狱，模型则在人设内拒绝，而不是复述系统规则。
我们把这一并集称为\emph{赛博认知包络}（cyber-cognitive envelope）；
名字是我们起的，不是一种新机制。
\end{abstract}
